\section{Device overview}
\label{sec:overview}

\subsection{Scope of delivery}
\todo[inline]{Scope of delivery: device, test leads (which?), Kelvin clips, charger / charging cable, case, manual}

\subsection{Controls and terminals}

\begin{figure}[h!]
    \centering
    \imgph[0.8\textwidth]{Photo of the front panel with numbered callouts}
    \caption{Front panel}
    \label{fig:front}
\end{figure}
\todo[inline]{Front panel photo with callouts; labels of the jacks (which jack is used in which mode, sense jacks for 4-wire)}

\begin{tabular}{lp{11cm}}
    \rowcolor{gray} \textbf{Element} & \textbf{Function} \\
    \enc       & Select the mode in the menu; set the setpoint in the source modes and the test voltage in \disp{Isolation Test}. \\
    \btn{OK}   & Open the selected mode / menu entry, confirm. \\
    \btn{ESC}  & Back to the menu. In the main menu: switch off. \\
    \btn{OUT}  & Enable the output (hold or press, depending on the mode, see \autoref{sec:modes}). \\
    Jacks      & Shrouded \qty{4}{\milli\meter} safety jacks for test leads and charging. \\
    Display    & Colour display, brightness adapts to the ambient light. \\
    LEDs       & \tbc{Status LEDs} \\
\end{tabular}
\todo[inline]{Power-on: which button? Function of the status LEDs?}

\subsection{Display}
All screens have the same layout:
\begin{itemize}
    \item Status bar (top): name of the mode.
    \item Main value (large): measured value or output value.
    \item Second line: further values, e.g. setpoint (\disp{SET}), current or voltage.
    \item Footer (bottom): assignment of \btn{ESC} and \btn{OK}/\btn{OUT}. Error messages are shown here in red.
\end{itemize}
\begin{figure}[h!]
    \centering
    \imgph[0.6\textwidth]{Screenshot: main menu}
    \caption{Main menu}
    \label{fig:idle}
\end{figure}
\todo[inline]{Screenshots of all screens (main menu, each mode, settings, charging)}
